\documentclass{article}
\usepackage{amsmath}
\title{Inline and Display Mathematics}
\author{ }
\date{}

\begin{document}
\maketitle

\section*{1. What is mathematical mode?}
LaTeX uses a special mode for writing mathematical expressions. There are two common ways to enter it: inline maths and display maths.

\section*{2. Inline mathematics}
Inline maths appears inside a sentence.

\textbf{LaTeX source code:}
\begin{verbatim}
The equation $a+b=c$ is an example.
\end{verbatim}

\textbf{Expected compiled output:}

The equation \(a+b=c\) is an example.

\textbf{Explanation:}
The dollar signs \verb|$...$| mark the beginning and end of inline maths. The equation stays on the same line as the sentence.

\section*{3. Display mathematics}
Display maths places an expression on its own line.

\textbf{LaTeX source code:}
\begin{verbatim}
\[
a+b=c
\]
\end{verbatim}

\textbf{Expected compiled output:}
\[
a+b=c
\]

\textbf{Explanation:}
The commands \verb|\[| and \verb|\]| begin and end an unnumbered display equation.

\section*{4. Practice}
Write \(x+y=z\) once inside a sentence and once on a separate line.
\end{document}
